\section{All integers, including supported zero, on the relative object}
\label{sec:all-integers-relative}
The local prime construction extends to the whole relative group
$D=\Q/\Z$. We keep its supported zero $0_D$ and the newly absent
element $\tau_D$ distinct. In particular the integer $e=0_\Z$ is
included in the scalar action, with its exact domain and codomain.

\subsection{Simultaneous arithmetic action and the original integral vector}
For $n\in\Z\setminus\{0\}$ the map $d\mapsto nd$ on $D$ is onto
with kernel $n^{-1}\Z/\Z$, of cardinality $|n|$. Therefore
\[
 (A_nf)(d)=f(nd),\qquad
 \|A_nf\|^2=|n|\|f\|^2,\qquad
 V_n=|n|^{-1/2}A_n
 \tag{IA1}
\]
on $\ell^2(D)$ with counting measure. These formulas follow by counting
each fibre, as in (PD7), and $A_mA_n=A_{mn}$ follows by composition.
The finite character pairings identify $\ell^2(D)$ with
$L^2(\widehat\Z)$, where $\widehat\Z=\prod_p\Z_p$ has Haar measure
one. Indeed every $d=a/N$ gives the character
$x\mapsto\exp(2\pi i a(x\bmod N)/N)$; the finite cyclic
orthogonality proof (PD5) applies with $N$ in place of $p^n$.
The functions factoring through the finite quotients form a dense
subspace, by compactness and the same approximation argument.

Let $\mathbb A_f$ denote the restricted product of the $\Q_p$ with
respect to $\Z_p$, with its original additive Haar measure
$\mu(\widehat\Z)=1$. For a positive rational $a$ define
\[
 (\mathcal U_aF)(x)=\sqrt a\,F(x/a),\qquad F\in L^2(\mathbb A_f).
 \tag{IA2}
\]
The finite-adelic modulus is $|a|_f=a^{-1}$: this follows first for
a prime from $|p|_p=p^{-1}$ and $|p|_q=1$ for $q\ne p$, and then
for all positive rationals by prime factorization. Changing variables
in (IA2) proves unitarity, and composition proves
$\mathcal U_a\mathcal U_b=\mathcal U_{ab}$.
For a positive integer $n$ its restriction to $L^2(\widehat\Z)$ is
the Fourier transform of $V_n$. The proof is the finite character sum
over the kernel of multiplication by $n$: it is zero off
$n\widehat\Z$, and equals $n$ times the original character on that
subgroup. This gives the formula
$\sqrt n\,\mathbf1_{n\widehat\Z}(x)F(x/n)$, including its factor.

Put $\xi=\mathbf1_{\widehat\Z}$, of norm one. For $a=m/n>0$ with
$\gcd(m,n)=1$,
\[
 \langle\mathcal U_a\xi,\xi\rangle=\frac1{\sqrt{mn}}.
 \tag{IA3}
\]
In fact $\widehat\Z\cap a\widehat\Z=m\widehat\Z$, as checked
at each valuation, and its measure is $1/m$. Multiplication by
$\sqrt{m/n}$ proves the displayed equality. In particular, for all
positive integers $m,n$,
\[
 \langle\mathcal U_m\xi,\mathcal U_n\xi\rangle
 =\frac{\gcd(m,n)}{\sqrt{mn}}.
 \tag{IA4}\label{eq:IA-gcd-Gram}
\]
This retains every integral scalar, not only primes or prime powers.
For a finite set of primes the joint scalar spectral measure is exactly
the product of the local measures (PD13): its joint Fourier coefficient
at $(k_p)$ is $\prod_pp^{-|k_p|/2}$ by (IA3). Expanding Laurent
polynomials proves the equality of their pairings with the actual
operators, just as in the proof of (PD13).

The determinant of the matrix (IA4) for $1\le m,n\le N$ is
\[
 \det\left(\frac{\gcd(m,n)}{\sqrt{mn}}\right)_{1\le m,n\le N}
 =\prod_{d=1}^N\frac{\varphi(d)}d>0.
 \tag{IA5}
\]
Here $\varphi(d)$ counts the units of $\Z/d\Z$, with $\varphi(1)=1$.
For completeness, sorting the integers $1,\ldots,k$ by the value of
$\gcd(a,k)$ proves $\sum_{d\mid k}\varphi(d)=k$: those with
$\gcd(a,k)=k/d$ correspond bijectively to residues coprime to $d$.
Thus, for the lower-triangular matrix $E_{md}=\mathbf1_{d\mid m}$,
\[
 (\gcd(m,n))_{m,n}=E\operatorname{diag}(\varphi(1),\ldots,\varphi(N))E^t.
\]
The diagonal of $E$ consists of ones, so its determinant is one.
Taking determinants and retaining the row and column factors
$m^{-1/2},n^{-1/2}$ proves (IA5).
This determinant identity is classical, associated with H. J. S. Smith
\cite{Smith1875}; the complete proof is given here. The calculation
specific to this construction is its realization as the Gram matrix
of the relative-prime operators on the original integral vector.

\subsection{Supported zero acts in an explicitly computed distribution space}
Let $\widetilde D=G(D)=D\sqcup\{\tau_D\}$ with its standard
split-semimodule scalar action. For supported $n\in\Z$,
$n\cdot d=nd$ is supported and $n\cdot\tau_D=\tau_D$.
For the absent scalar $\tau$, every product is $\tau_D$.
Let
\[
 \mathcal T=\C^{(D)}\oplus\C\delta_{\tau_D},\qquad
 \mathcal T'=\C^D\oplus\C,
 \qquad \langle f,h\rangle=\sum_{d\in D}f(d)h(d)+c_f c_h.
 \tag{IA6}
\]
The first factor of $\mathcal T$ consists of finite-support functions,
so the displayed pairing is defined for every $h\in\mathcal T'$.
It is the full algebraic dual pairing. No square-summability condition
is imposed on a distribution in $\mathcal T'$.

Pullback by the scalar maps defines operators $R_a$ on all functions
$\mathcal T'$ satisfying $R_aR_b=R_{ab}$ and $R_1=I$. For nonzero
supported $n$, $R_n$ restricts to (IA1) on $D$ and fixes the absent
coordinate. For the two distinct zeros it gives, exactly,
\[
 R_e(f,c)=(f(0_D)\mathbf1_D,c),\qquad
 R_\tau(f,c)=(c\mathbf1_D,c).
 \tag{IA7}\label{eq:IA-two-zero-operators}
\]
In particular
\[
 R_e\delta_{0_D}=(\mathbf1_D,0),\quad
 R_\tau\delta_{0_D}=0,\quad
 R_e\delta_{\tau_D}=\delta_{\tau_D},\quad
 R_\tau\delta_{\tau_D}=(\mathbf1_D,1).
 \tag{IA8}
\]
Every equality follows by evaluating at a supported $d$ and at
$\tau_D$. They prove, rather than identify away, the difference
between the action of the integer zero and the absent scalar.
Composition of the original scalar maps proves the representation law
also when $a$ or $b$ is either zero. This is a representation of the
\emph{multiplicative scalar monoid}; no false additive identity
$R_{a+b}=R_a+R_b$ is asserted for pullbacks of functions.

There is an exact reason to keep the distribution space (IA6).
The nonzero function $\mathbf1_D$ has infinite counting-measure norm.
Thus $R_e$ has no bounded extension on $\ell^2(D)$ agreeing with
(IA8). Its full action nonetheless exists on $\mathcal T'$ and its
restriction $\mathcal T\to\mathcal T'$ is given by (IA7).
Its Fourier transform on the arithmetic component is
\[
 f(0_D)\mathbf1_D
 \longleftrightarrow
 \left(\int_{\widehat\Z}\mathcal F f(x)\,dx\right)\delta_0.
 \tag{IA9}\label{eq:IA-zero-Fourier}
\]
To prove this without a convergence claim, test against a character
polynomial $F=\sum_da_d\chi_d$. The distribution with all Fourier
coefficients one takes the value $\sum_da_d=F(0)$, which is exactly
$\delta_0$. Orthogonality (PD5) gives
$f(0_D)=\int\mathcal F f$. This proves both sides of (IA9) on every
test. The algebraic transpose of $R_e$ on distributions is also
represented on finite-support tests: it sends
$(f,c)$ to $((\sum_df(d))\delta_{0_D},c)$. Direct substitution in
(IA6) verifies the transpose identity. Thus the mass and supported-zero
evaluation terms are explicit adjoint boundary maps.

The integer $0$ has therefore remained in every scalar map. Its action
lands at a boundary distribution that the Hilbert completion alone
does not contain. The maps (IA7)--(IA9) specify this enlargement and
its relation to the positive Hilbert operators (IA1)--(IA4). Any
purity or duality argument using that Hilbert realization must retain
this computed zero-scalar boundary and its transpose, rather than
quietly delete the integer zero from the programme.
